\section{Background}
\label{sec:background}

This section exists so the rest of the report can be read without looking
anything up. A reader who already knows WCA notation and CTC can skip to
\cref{sec:data}.

\subsection{The cube as a group}

Label the six faces \wca{U}p, \wca{D}own, \wca{L}eft, \wca{R}ight,
\wca{F}ront, \wca{B}ack. A \emph{quarter turn} rotates one face by
$90^\circ$: \wca{R} clockwise (looking at that face), \wca{R'}
anticlockwise. \wca{R2} is a half turn and in this system is always
represented as two quarter turns, which matters more than it sounds
(\cref{sec:qtm}).

Configurations of the cube form a group $G$ under composition of moves,
with $|G| = 43{,}252{,}003{,}274{,}489{,}856{,}000 \approx 4.3\times10^{19}$.
The concrete representation used throughout the code is the standard
\emph{cubie} encoding: a permutation of 8 corners with orientations in
$\mathbb{Z}_3$, and a permutation of 12 edges with orientations in
$\mathbb{Z}_2$, i.e.\ a vector in
$S_8 \times \mathbb{Z}_3^8 \times S_{12} \times \mathbb{Z}_2^{12}$ subject
to three parity constraints. Composition is array indexing, so a beam of
thousands of hypotheses can be advanced by one move with a single vectorised
operation --- this is why the decoder in \cref{sec:reconstruct} is
affordable at all.

Two facts about $G$ are load-bearing later.

\begin{description}[leftmargin=1.2em,style=nextline,itemsep=2pt]
\item[Every quarter turn flips permutation parity.] A quarter turn is a
4-cycle on corners and a 4-cycle on edges, each odd, so the total
permutation parity flips. Consequently the parity of the number of quarter
turns in a word is determined by its endpoints. A hypothesis whose remaining
move budget has the wrong parity provably needs an extra insertion or
deletion, which gives the decoder a free admissible lower bound
(\cref{sec:heuristic}).
\item[God's number is 20 in the half-turn metric and 26 in the quarter-turn
metric.] \label{sec:qtm} Every state is at most 20 half turns, or 26 quarter
turns, from solved \cite{rokicki2014,rokicki2010}. Because this system's
alphabet has no half-turn symbol, \emph{all} of its counting is
quarter-turn, and a threshold calibrated on the more famous number 20 would
be set below the quantity it is supposed to bound. \cref{sec:count-gate}
depends on getting this right.
\end{description}

\subsection{How humans actually solve, and why it matters}

The dominant method, CFOP, proceeds in four phases: a cross on one face
($\sim$8 moves), four \emph{first-two-layers} (F2L) pairs inserted one at a
time ($\sim$30 moves), then orientation and permutation of the last layer
(OLL/PLL) executed as memorised \emph{algorithms} --- fixed 8--15 move words
drawn from a personal repertoire of a few dozen.

Two consequences run through this report. First, a human solve is far longer
than an optimal one: 50--130 quarter turns against a solver's $\le 26$. That
gap is the entire basis of the anticheat count gate (\cref{sec:count-gate}).
Second, the last third of a solve is drawn from a small closed vocabulary,
which suggests a language model or algorithm prior should help a lot. The
repository pursued that idea hard; \cref{sec:negatives} records that it
bought $0.2$ points.

\subsection{Connectionist temporal classification}
\label{sec:ctc-bg}

CTC \cite{graves2006} is the standard answer to ``I have a sequence of
$T$ frames and a target sequence of $U \ll T$ symbols, and I do not know
which frame produced which symbol.''

Let $\mathcal{A}$ be the label alphabet ($|\mathcal{A}| = 12$ here) and let
$\mathcal{A}' = \mathcal{A}\cup\{\varnothing\}$ add a \emph{blank}. The
network emits, for each frame $t$, a distribution $y^t$ over $\mathcal{A}'$
(so 13 columns). A \emph{path} $\pi \in \mathcal{A}'^T$ assigns one symbol
to every frame, with probability
\begin{equation}
p(\pi \mid x) \;=\; \prod_{t=1}^{T} y^t_{\pi_t},
\label{eq:pathprob}
\end{equation}
which assumes conditional independence across frames given the network
output --- an assumption we return to in \cref{sec:ctc-limits}.

The \emph{collapse} map $\mathcal{B}$ takes a path to a label sequence by
first merging runs of identical symbols and then deleting blanks:
$\mathcal{B}(\wca{R}\,\wca{R}\,\varnothing\,\wca{R}\,\wca{U}) =
\wca{R}\,\wca{R}\,\wca{U}$. The probability of a label sequence $\ell$ is
the total probability of every path that collapses to it:
\begin{equation}
p(\ell \mid x) \;=\; \sum_{\pi \in \mathcal{B}^{-1}(\ell)} p(\pi \mid x),
\label{eq:ctcmarginal}
\end{equation}
and training minimises $-\log p(\ell^\star \mid x)$ for the true label
sequence $\ell^\star$. The sum in \cref{eq:ctcmarginal} has exponentially
many terms and is computed in $O(TU)$ by a forward--backward recursion over
an extended target $\ell' = \varnothing\, \ell_1\, \varnothing\, \ell_2
\cdots \ell_U\, \varnothing$ of length $2U+1$: with
$\alpha_t(s)$ the total probability of all prefixes of paths that have
consumed $\ell'_{1:s}$ by frame $t$,
\begin{equation}
\alpha_t(s) \;=\; y^t_{\ell'_s}\cdot
\begin{cases}
\alpha_{t-1}(s) + \alpha_{t-1}(s-1)
  & \text{if } \ell'_s = \varnothing \text{ or } \ell'_s = \ell'_{s-2},\\[2pt]
\alpha_{t-1}(s) + \alpha_{t-1}(s-1) + \alpha_{t-1}(s-2)
  & \text{otherwise.}
\end{cases}
\label{eq:ctcforward}
\end{equation}
The case split is the entire content of CTC and is worth internalising: the
three-term recursion is the ordinary ``advance to the next target symbol''
step, and the two-term case says that you may \emph{not} skip a blank when
the symbol before it is the same as the one after it. That restriction is
what makes a genuine repeat expressible --- and it is also the source of
CTC's one structural failure mode here, discussed next.

\paragraph{The structural miss.} Two identical consecutive labels require an
intervening blank frame. So a true \wca{R}\,\wca{R} whose two turns are one
frame apart at 30\,fps \emph{cannot} be emitted as two symbols, no matter how
good the network is. In this corpus that costs
\CTCFloorHoldout\% of held-out onsets, i.e.\ a ceiling of
\CTCCeilingHoldout\% per-move accuracy that no amount of training removes.
It is a real cost of the design, and it is quantified rather than waved at
because it bounds every result in \cref{sec:experiments}. Half turns are
\emph{not} the mechanism (a half turn's two quarter turns are a median
$\sim$0.33\,s apart, ten frames, so they are cleanly separable); the
mechanism is genuine two-handed simultaneity.

\subsection{Related problem shapes}

The pipeline is an assembly of three well-studied shapes, and naming them
correctly is useful for knowing what literature to steal from.

\begin{itemize}[leftmargin=1.3em,itemsep=2pt]
\item \textbf{Temporal action spotting / onset detection.} Finding the
instant of a brief event in continuous video. The pre-CTC arm of this
project was exactly this: a per-frame score with Gaussian targets, peak
picked. The music-onset-detection literature \cite{bock2012} uses the same
construction, including the Gaussian-smeared targets and the
minimum-separation peak picker, and the failure mode observed here (two
events inside one hump) is the same one reported there.
\item \textbf{Sequence transduction without alignment.} CTC \cite{graves2006}
came from speech; the version used here is unmodified, with a
character-level analogue in which ``characters'' are quarter turns. The
prefix beam search of \cref{sec:prefixbeam} is Hannun et al.'s
\cite{hannun2014}.
\item \textbf{Noisy-channel decoding against a hard constraint.} Given a
corrupted string and a formal constraint the true string must satisfy, find
the minimum-cost correction. This is spell correction, it is constrained
decoding in MT, and here the constraint is membership of a specific coset of
the cube group. \cref{sec:reconstruct} is a beam search over that.
\end{itemize}

\noindent The one thing that has \emph{no} obvious literature analogue is
the verification framing: a decoder whose output is used as evidence about
whether the input was honest. The closest analogue is presentation-attack
detection in biometrics, and the borrowed lesson --- that a detector
evaluated only on genuine inputs has not been evaluated --- is why
\cref{sec:falsifiability} exists.
